% IEEE Access — figure/table captions

% fig01_pipeline_overview: End-to-end TDE pipeline from ESRF CT acquisition to labeled synthetic deformation dataset. Stages correspond to the three processing modules: data preparation, 3-D projection, and synthetic deformation generation.
% table01_pipeline_stages: TDE pipeline stages, responsible parties, inputs, outputs, and key functions.
% fig02_module_dependency: Software module dependency graph of the TDE pipeline. Nodes represent Python modules; directed edges denote import or data-flow dependencies.
% fig03_data_flow: Data flow diagram of the TDE pipeline showing transformations applied at each processing step and the intermediate data formats.
% fig04_specimen_distribution: Number of 2-D projection images generated per fossil specimen. Each specimen contributes exactly 30 projections (6 orthographic + 24 rotational), yielding a balanced dataset.
% fig05_dataset_composition: Dataset composition showing equal contribution of all eight archosaur specimens to the TDE synthetic deformation dataset. Each specimen contributes 120 synthetic images (12.5%).
% fig06_dataset_growth: Dataset scale at each stage of the TDE pipeline, demonstrating the multiplicative effect of multi-view projection (×30) and synthetic deformation (×4) on the base specimen count.
% fig07_projection_types: Breakdown of the 30 projection views generated per specimen: 6 orthographic views (3 axes × 2 directions) and 24 rotational views (0° to 345° at 15° increments) of the sagittal MIP.
% fig08_deformation_distribution: Distribution of synthetic images across four taphonomic deformation categories. All categories are balanced with 240 images each, producing a class-balanced training dataset.
% fig09_projection_pipeline: Streaming Maximum Intensity Projection (MIP) architecture. JP2 slices are decoded one at a time; three orthogonal MIPs are accumulated in a single pass, reducing peak memory from O(D \cdot H^2) to O(H^2).
% fig10_sample_projections: Representative 2-D Maximum Intensity Projection images from four archosaur specimens. Top row: orthographic coronal views. Bottom row: rotational views (45°, 90°, 135°, 180°). All images are 224×224 grayscale PNG.
% fig11_all_specimens_grid: Complete gallery of orthographic coronal (ortho\_00) Maximum Intensity Projection images for all eight archosaur specimens in the TDE dataset. Images are 224×224 grayscale PNG.
% fig12_voxel_sizes: CT voxel resolution (µm) for each archosaur specimen. Voxel sizes span two orders of magnitude (1.28–107.16 µm), reflecting the diverse scan resolutions available in the ESRF Paleo CT dataset.
% fig13_transformation_matrices: Transformation matrix representations for the four taphonomic deformations. (a) Compression: $c \in [0.3, 0.9]$. (b) Shearing: $k \in [-0.5, 0.5]$. (c) Stretching: $s_x, s_y \in [1.0, 1.5]$. (d) Dissolution: intensity subtraction via ellipse mask.
% fig14_deformation_geometry: Geometric illustration of the four taphonomic deformation types applied to a unit square. The dashed square shows the original boundary; the solid polygon shows the deformed result.
% fig15_deformation_gallery: Deformation gallery for the araripesaurus specimen (ortho\_00 view). Each column shows: original projection, deformed image, and binary deformation mask. Rows correspond to compression, shearing, stretching, and dissolution respectively.
% fig16_parameter_distributions: Parameter distributions for all four deformation types across the full TDE dataset (n=240 per type, seed=42). Parameters are uniformly sampled: $c \sim \mathcal{U}(0.3,0.9)$, $k \sim \mathcal{U}(-0.5,0.5)$, $s_x, s_y \sim \mathcal{U}(1.0,1.5)$, $\lambda \sim \mathcal{U}(0.3,0.8)$.
% fig17_stretching_scatter: Joint distribution of stretching parameters $s_x$ and $s_y$ (n=240). Both are independently sampled from $\mathcal{U}(1.0,1.5)$, confirming near-zero correlation and independent axis scaling.
% fig18_compression_progression: Effect of increasing compression factor $c$ on specimen appearance. A lower value of $c$ produces stronger vertical compression, simulating sedimentary overburden deformation.
% fig19_dissolution_masks: Example binary deformation masks for the dissolution deformation type. Each mask encodes the elliptical region where intensity was attenuated. Masks are provided for Grad-CAM explainability validation.
% fig20_mask_coverage_stats: Distribution of deformation mask coverage (\% of image pixels changed) per deformation type. Dissolution masks reflect the random ellipse area; compression, shearing, and stretching masks reflect the affine-transformed region.
% fig21_intensity_distributions: Pixel intensity distributions (0–255) for original projection images and each synthetic deformation type, sampled from araripesaurus\_\_BSPG-1982-I-90. Dissolution is the only deformation that shifts the mean intensity downward.
% fig22_folder_structure: Canonical folder hierarchy of the TDE dataset. Each specimen directory contains an ESRF-standard \texttt{org\_slices/} directory with JP2 CT slices and a \texttt{mesh/} directory for 3-D surface meshes.
% fig23_function_call_graph: Key function call graph for the TDE projection and deformation pipelines. Edge direction indicates caller → callee. Nodes are colored by processing module.
% fig24_4deformation_panel: Publication-quality comparison panel of the four taphonomic deformations applied to the Archaeopteryx\_london specimen (BMNH-37001, ortho\_00). Each row shows original projection, synthetic deformed image, and pixel-wise binary mask.
% fig25_rotational_strip: Twelve representative rotational MIP views of the halszkaraptor specimen at 30° increments. Views are generated by rotating the sagittal MIP (axis-1 projection) in 2-D after accumulation.
% fig26_mip_algorithm: Conceptual illustration of the streaming Maximum Intensity Projection (MIP) algorithm. JP2 slices are processed one at a time; a running \texttt{np.maximum} accumulates the global MIP in O(H·W) memory regardless of volume depth D.
% fig27_percentile_normalization: Percentile normalization (1--99\%) applied to a raw MIP accumulator output from the araripesaurus specimen. Left: raw uint16 histogram; Right: normalized uint8 output after clipping and rescaling.
% fig28_taxonomy_voxel: Archosaur specimen taxonomy and CT voxel resolution. Marker size encodes voxel resolution (µm). The dataset spans birds (Aves), non-avian dinosaurs (Dromaeosauridae), and flying reptiles (Pterosauria).
% fig29_dataset_summary: TDE dataset summary statistics. The dataset contains 240 projection images and 960 synthetic images with corresponding binary masks and per-image labels across four balanced deformation classes.
% fig30_csv_schema: Schema of the \texttt{synthetic\_labels.csv} output file. Each row describes one synthetic image with its specimen, deformation type, sampled parameters, output filenames, and corresponding binary mask filename.
